% ====================================================================
% 서지 (15 외부 문헌 + 내부자료 numverif2026 = 16 bibitem)
% ====================================================================
\begin{thebibliography}{99}
\bibitem{bernardi1985} D. Bernardi, E. Pawlikowski, J. Newman, ``A General Energy Balance for Battery Systems,'' \emph{J. Electrochem. Soc.} \textbf{132}(1), 5--12 (1985). DOI: 10.1149/1.2113792.
\bibitem{newman} J. Newman, K. E. Thomas-Alyea, \emph{Electrochemical Systems}, 3rd ed., Wiley (2004). (OCV($T$): slope $=\Delta S$, intercept $=\Delta H$; lattice-gas insertion thermodynamics.)
\bibitem{huggins2009} R. A. Huggins, \emph{Advanced Batteries: Materials Science Aspects}, Springer (2009). (격자기체 삽입 전극 열역학·Bragg--Williams 평형 전위.)
\bibitem{bazant2013} M. Z. Bazant, ``Theory of Chemical Kinetics and Charge Transfer Based on Nonequilibrium Thermodynamics,'' \emph{Acc. Chem. Res.} \textbf{46}(5), 1144--1160 (2013). DOI: 10.1021/ar300145c. (regular-solution lattice-gas OCV·상분리 임계 $\Omega=2RT$.)
\bibitem{dahn1991} J. R. Dahn, ``Phase diagram of Li$_x$C$_6$,'' \emph{Phys. Rev. B} \textbf{44}, 9170 (1991). DOI: 10.1103/PhysRevB.44.9170.
\bibitem{ohzuku1993} T. Ohzuku, Y. Iwakoshi, K. Sawai, ``Formation of Lithium-Graphite Intercalation Compounds in Nonaqueous Electrolytes and Their Application as a Negative Electrode for a Lithium Ion (Shuttlecock) Cell,'' \emph{J. Electrochem. Soc.} \textbf{140}, 2490 (1993). DOI: 10.1149/1.2220849.
\bibitem{reynier2003} Y. Reynier, R. Yazami, B. Fultz, ``The entropy and enthalpy of lithium intercalation into graphite,'' \emph{J. Power Sources} \textbf{119--121}, 850--855 (2003). DOI: 10.1016/S0378-7753(03)00285-4.
\bibitem{allart2018} D. Allart \emph{et al.}, ``Model of Lithium Intercalation into Graphite by Potentiometric Analysis with Equilibrium and Entropy Change Curves of the Graphite Electrode,'' \emph{J. Electrochem. Soc.} \textbf{165}(2), A380 (2018). DOI: 10.1149/2.1251802jes.
\bibitem{occupation2019} M. P. Mercer, M. Otero, M. Ferrer-Huerta, A. Sigal, D. E. Barraco, H. E. Hoster, E. P. M. Leiva, ``Transitions of lithium occupation in graphite: A physically informed model in the dilute lithium occupation limit supported by electrochemical and thermodynamic measurements,'' \emph{Electrochimica Acta} \textbf{324}, 134774 (2019). DOI: 10.1016/j.electacta.2019.134774. [방법 수준 인용 --- 원문 식 미확보.]
\bibitem{chemmater2015} S. Konar, U. Häussermann, G. Svensson, ``Intercalation Compounds from LiH and Graphite: Relative Stability of Metastable Stages\ldots,'' \emph{Chem. Mater.} \textbf{27}(7), 2566--2575 (2015). DOI: 10.1021/acs.chemmater.5b00235. [형성 $\Delta H$ calorimetry --- abstract tier.]
\bibitem{jpcc2021} J. Haruyama, S. Takagi, K. Shimoda, I. Watanabe, K. Sodeyama, T. Ikeshoji, M. Otani, ``Thermodynamic Analysis of Li-Intercalated Graphite by First-Principles Calculations with Vibrational and Configurational Contributions,'' \emph{J. Phys. Chem. C} \textbf{125}(51), 27891--27900 (2021). DOI: 10.1021/acs.jpcc.1c08992. [abstract tier.]
\bibitem{msmr_partI} T. R. Garrick, B. J. Koch, M. Choi, X. Du, A. M. Adeyinka, J. A. Staser, S.-Y. Choe, ``Quantifying the Entropy and Enthalpy of Insertion Materials for Battery Applications Via the Multi-Species, Multi-Reaction Model,'' \emph{J. Electrochem. Soc.} \textbf{171}(2), 023502 (2024). DOI: 10.1149/1945-7111/ad1d27. [★Chapter 1 이 인용하는 ECS Advances 의 ``Part 1''(MCMB 흑연 파라미터화, DOI: 10.1149/2754-2734/ad7d1c)과는 \emph{별개 논문} --- 두 문헌이 공유하는 것은 Part II(DOI: 10.1149/1945-7111/ad70d9) 뿐이다.]
\bibitem{msmr_partII} A. Paul, K. Wolfe, M. W. Verbrugge, B. J. Koch, J. S. Lowe, J. Trembly, J. A. Staser, T. R. Garrick, ``Quantifying the Temperature Dependence of the Multi-Species, Multi-Reaction Model: Part~II. Estimation of Entropy Coefficient for Meso-Carbon Micro-Bead Graphite,'' \emph{J. Electrochem. Soc.} \textbf{171}(10), 103505 (2024). DOI: 10.1149/1945-7111/ad70d9.
\bibitem{standardised2024} A. Hales, J. Bulman, ``A Standardised Potentiometric Method for the Effective Parameterisation of Reversible Heating in a Lithium-Ion Cell,'' \emph{J. Electrochem. Soc.} \textbf{171}(5), 050535 (2024). DOI: 10.1149/1945-7111/ad4918.
\bibitem{hysteresis2018} I. Zilberman, A. Rheinfeld, A. Jossen, ``Uncertainties in entropy due to temperature path dependent voltage hysteresis in Li-ion cells,'' \emph{J. Power Sources} \textbf{395}, 179--184 (2018). DOI: 10.1016/j.jpowsour.2018.05.052.
\bibitem{numverif2026} 본 연구 내부 수치 검증(Derivative A), \code{Anode\_Fit\_v1.0.19} 4-전이 흑연 staging 파라미터, 유한차분 vs.\ 가중식 비교 (2026). [내부 자료 --- 표시 정밀도 일치 확인, 본문 \S\ref{ssec:overlap}$\cdot$\S\ref{ssec:blend} 요지.]
\end{thebibliography}

% 자산: [C-120] bernardi1985(eq:qrev 원출처) · [C-121] newman(OCV(T) slope=ΔS) · [C-122] huggins2009(격자기체·BW) ·
%   [C-123] bazant2013(regular-solution·Ω=2RT) · [C-124] reynier2003(vib 음baseline) · [C-125] allart2018(정량값·tab:ds) ·
%   [C-126] occupation2019([방법수준]dilute 발산) · [C-127] chemmater2015([형성ΔH abstract]) · [C-128] jpcc2021([abstract]제일원리) ·
%   [C-129] ★msmr_partI(Garrick 2024, Ch1 Part1과 별개·식별함정) · [C-130] msmr_partII(Paul 2024·MCMB) ·
%   [C-131] standardised2024(calorimetry) · [C-132] hysteresis2018(경로의존) · [C-133] ★numverif2026([내부자료]파생A 원천)
